\scanpage{228}
cas et dans ce cas seulement, on peut considérer
que $\widehat Z=G_{\varphi_{\unsure{0|\varphi}}}$ opère sur $\Pi_Z$ (et non seulement sur $\varphi_0$),
-- et alors $\Pi_Z=\Pi_{\widehat Z}$ est un torseur sous $\widehat Z$.

Nous nous intéressons, pour $(\Gamma,Z\subset A_\Gamma=\mathrm{Aut}(\Gamma))$
donné, avec $\Gamma=(S,S^\vee,R)$, $Z$ ss-groupe commutatif,
\deleted{aux} \inserted{à} la donnée d'un sous-torseur $\pi$ de $\Pi_Z$,
sous l'opération de $\varphi$. Dans le cas d'une
$\Gamma$ biséparée, \unsure{cel|cela} signifie qu'on s'intéresse aux
``multigraphes séparés à groupe d'homothéties $Z$.''
\deleted{\unsure{On}}
\[
 \underset{\substack{\text{sous}\\Z\text{-torseur}}}{\pi}\subset\Pi_Z=\text{ens des } Z\text{-jauges bijectives}.
 \tag*{(30)}
\]
On trouve alors des bijections canoniques
\[
 S^\vee\simeq\pi\wedge_Z(S!),\qquad S\simeq\pi^\vee\wedge_Z(S^\vee!)
 \tag*{(31)}
\]
dont la \unsure{deuxième|dernière} peut être considérée, soit comme
déduite de la première, soit comme étant
l'isomorphisme du $=$ type que la première,
relatif à la situation duale
\[
 \left(\underset{\substack{\Vert\\(S^\vee,S,{}^sR)}}{\Gamma'},
 \underset{\substack{\Vert\\ {}^sZ=\{\{\lambda_{S^\vee},\lambda_S\}\mid\lambda\in Z\}}}{Z'},
 \underset{=\{\varphi^{-1}\mid\varphi\in\pi\}}{\pi^{\deleted{\unsure{\vee}}'}\ \deleted{\unsure{Z'}}}\right)
 \tag*{(32)}
\]
en notant que
\[
 \varphi\longmapsto\varphi^{-1}\qquad \pi\longrightarrow\pi'
 \tag*{\deleted{\unsure{(33)}}}
\]
\deleted{\unsure{est un}} \inserted{satisfait à}
\[
 (T_\lambda\varphi)^{-1}=T'_{\lambda^{-1}}(\varphi^{-1})
\]
et définit donc un isom can. de $Z$-torseurs
\[
\begin{aligned}
 \pi'\simeq\pi^{[-1]}\overset{\mathrm{déf}}={}&\pi\text{ avec opn de }Z\text{ opposée de}\\
 &\text{l'opération donnée } T_\lambda,\text{ i.e.}\\
 &\lambda\text{ opérant par }T_{\lambda^{-1}}\ (\overset{\mathrm{déf}}=T'_\lambda).
\end{aligned}\tag*{(33)}
\]
% archive sheet number: 227

